\section{Meditations on the Good Life}

\begin{quotex}
Death, when understood by men, is immortality; but not being understood by the ignorant it is death. This death should not be feared; what should be feared is the perdition of the soul, which is ignorance of God; for this is terrible for the soul. \flright{\textsc{St Anthony the Great}, \emph{Counsels on the Character of Men and on the Virtuous Life}}

In tantum Deus cognoscitur, in quantum amatur. \flright{\textsc{St Bernard}}

(God is known to the extent He is loved.) 

\end{quotex}
In recent days we have been advised on how to be good:

\begin{itemize}
\item By the president: Support the right to traffic in human body parts 
\item By a candidate: Make America great again 
\item By the pope: Turn off the air conditioner 
\end{itemize}
In \emph{The Sparkling Stone}, \textbf{Blessed John of Ruysbroeck} provides a more traditional perspective, as summarized in the following outline:

\begin{enumerate}
\item \textbf{How a man becomes good} 

\begin{enumerate}
\item Clean conscience 
\item Obedience to God 
\item The glory of God 
\end{enumerate}
\item H\textbf{ow a man becomes inward} 

\begin{enumerate}
\item A heart unencumbered with images 
\item Spiritual freedom in his desires 
\item The feeling of inward union with God 
\end{enumerate}
\item \textbf{How a man can see God} 

\begin{enumerate}
\item The foundation of his being is abysmal 
\item His inward exercise is wayless 
\item His indwelling is a divine fruition 
\end{enumerate}
\end{enumerate}
\paragraph{Becoming Good}
First a man must have a clean conscience that is free of sin. There are three elements of sin.

\begin{enumerate}
\item It must be a serious matter 
\item Knowledge that it is a serious matter 
\item Performed willfully 
\end{enumerate}
The pagan view found, for example, both among the ancient Greeks and Hindus, stops at the second element. For them, ignorance, or avidya, alone would lead to the clean conscience. Then the practice of the virtues would follow from that knowledge. However, this didn't take into account that, besides the ignorance of the mind, there is also a weakness of the will.

Hence, knowledge alone is insufficient, one must consciously follow God. John also adds that a man must be ``obedient to his own proper convictions." Of course, the most proper convictions must be based on true knowledge. At a higher stage, John claims that even knowledge must be left behind in a state of learned ignorance, in which a man becomes free of such sin.

Finally, a good man should act for the glory of God. John explains:

\begin{quotex}
If it happens that by reason of his business or the multiplicity of his works, he has not always God before his eyes, at least there should be established in him the intention and desire to live according to the dearest will of God. 

\end{quotex}
This three qualities make a man good, so he is prepared for the next stages.

\paragraph{Becoming Inward}
John advises that anyone who considers himself spiritual needs to observe himself. Images in the imagination come from the flesh, not from God who is Spirit so no images are possible of Him. According to an old Norse adage,

\emph{The divine sleeps in the rock, breathes in the plant, dreams in the animal, and wakes in man.}

This is related to Plotinus' path from the psyche, through nous, to God. The material world is both physical (the ``rock") and psychical (the ``animal"). The animal soul is material and is full of images. Thus it lives in a dream-like state, not yet self-conscious. To the extent that a man, who does indeed possess an animal soul, lives the same way, he is not yet awake and not living a fully human life. So for the heart to become unencumbered with images, it must rise to the level of the nous by consciously observing himself.

As a man becomes free of worldly images, he becomes inwardly free. He becomes oriented toward God rather than toward the world. As imageless and free, he beings to feel a true union with God within; this perfects his inward and spiritual life. This acts like a feedback loop. The feeling of union, the desiring power of the soul (\emph{epithetikos} or \emph{eros}) becomes enticed to seek deeper union. God's knowledge of you is your knowledge of God; your willing of your own being is also God's willing.

\paragraph{Seeing God}
This union is experienced as ``abysmal", i.e., arising from the measureless depths, unbounded heights and an endless horizon. He plunges into the depths of the soul and ascends to the heights. The horizon of Being always recedes. Our knowledge of God is never complete so we dwell in a ``knowledge that is ignorance". By dying to things, the soul feels itself to be one life with God.

The next point is above reason and without condition. The inward drawing of the Divine Unity is experienced according to the measure of his love and the manner of his spiritual, interior exercises. Spiritually naked and unencumbered by images, he finds an Eternal Light in the inmost part of his spirit. This Light reveals the ``eternal demand of the Divine Unity". The more he feels it, the more he craves it.

The fruition is the loving adherence to God. There is an inner silence, so that thoughts and images begin to lose their impact and compulsion. The spirit then beholds a Darkness, where discursive thought cannot penetrate. Hence, the spiritual exercise is ``wayless", that is, there is no straight path to that Darkness, no mechanical process, no repetition of mantras. Rather one must be drawn into it, first though interior silence and the will to transcend discursive thought to a direct intuition.



\flrightit{Posted on 2015-08-03 by Cologero }

\begin{center}* * *\end{center}

\begin{footnotesize}\begin{sffamily}



\texttt{Mark Citadel on 2015-08-06 at 12:48 said: }

The contrast between John's view of the good could not stand in more stark contrast to the pablum we hear from the Modernist wardens could it?

``He becomes oriented toward God rather than toward the world."

Society follows the same metric. Oriented towards telluric interests it will be evil, but oriented up towards the divine, it will find its existence blessed, if only by natural forces and factors that govern success and decline.


\hfill

\texttt{Cologero on 2015-08-07 at 08:39 said: }

Mark, another way to put it is that what the world finds lovable and what God considers lovable are quite different.


\hfill

\texttt{Aleksandar Gueric on 2015-10-02 at 08:34 said: }

About becoming inward, Novalis wrote that "mysterious path leads within" and that "In us, or nowhere, is eternity with its worlds — the past and the future." This quote, for me,contains an explanation of a peculiar feeling, or sensation, I have had in me since very young age.

Past, for me, was never something gone but a living experience reachable through the mere power of concentration while objects around us – trees and rocks, when absorbed with intiutive knowing become living parts of the self. So, tree didn't decided to speak to me, for example, when I walked by it on one autumn day when I suddenly felt it's perspective – it's experience of time and change of seasons quite different from every day human one. It was only an intuition or deep knowledge and understanding of things behind their material forms and natualistic mechanics.. I think that simultaneous experience of human consciousness and totality of all other non human consciousness as well as awareness of timelesness through the simple focus of mind, must be what people call "being awake".

Novalis' Magical Idealism sparked this kind of imagination and intuition ever since reading about him as a child in a big encyclopedia we had at home. I recognized his language although I had no intellectual capacity to know why, only intuitively how, through a symbolic collage of metalanguage. Many years forward, Lama Anagarika Govinda spoke with the same language to me in "Foundations of Tibetan Mysticism", book that have found me by pure accident. Then I realized he quoted Novalis a lot and, at one point, even considered being his reincarnation. I think that he was close. They both spoke the same language, walked the esotericaly same path, perhaps. This language, or knowledge, I date back to Plato who formulated and crystalized it's essence. He was definitely an Initiate.

My peregrination during the turbulent 20s and falling into one trap after another in my constant search for something, which led me across 3 continents, finally seemed to come to an end with the discovery of another author who spoke the secret language I understood, and who also loved to quote Novalis, my initial sparkle – Baron Julius Evola. Gate has been opened enabling me to finally move forward transcending in the spiral way out. Novali's spiral path inwards.

This path is polarized or centered around the notion of a sacred center, which was, for me, always in Hyperborea. Real, unreal, mythical, physical, etheric, poetic, all in one. Magnetism was so strong that I envisioned myself bathing in Aurora Borealis, a cosmic North, source of light and purity. This is the greatest Ideal for me and true Home, which is "everywhere and nowhere".

I apologize for what must have been a crazy post. This site carries the same message and understanding which guided me through years of my journey and I'm grateful to find it.


\end{sffamily}\end{footnotesize}
